\chapter{货币是怎么"凭空"出现的}
\chaptermeta{07}

\begin{ledetext}
第 2 章说过，你的存款不在银行的金库里。这一章要说的更难听一点：那笔钱在被借给你之前，压根不存在。
\end{ledetext}

\section{2024 年 6 月的一个下午，300 万从空气里长出来}

陈宁 34 岁，在杭州一家做工业相机的公司当结构工程师，月薪 2 万 3，到手 1 万 8 出头。2024 年 6 月他在余杭买了套 486 万的二手房，卖家是个想换大三居的中学老师，姓项。首付 186 万，剩下 300 万找银行贷。

贷款金额之所以正好是 300 万这么齐整的数，不是他挑的，是评估价乘以贷款成数的上限卡在那儿，他顶格贷满了。

那天下午三点多，客户经理发来一条微信：过了。

把这一秒放慢。批贷那一刻，银行到底干了什么？

多数人脑子里是这幅画面：银行有个大池子，装着王阿姨、你妈、隔壁老陈存进去的钱。信贷员拿个瓢舀出 300 万倒进陈宁账户，水位下降 300 万。

这幅画很直观，很符合常识，而且\textbf{是错的}。

真实发生的事情是，一位系统操作员在核心账务系统里敲下了两行记录。

\begin{bsheet}
  \bsside{放贷前 · 某支行（万元）}{%
    \bsrow{准备金}{47}
    \bsrow{持有的债券}{63}
    \bsrowtotal{资产合计}{110}
  }
  \bsside{负债 + 净值}{%
    \bsrow{客户存款}{98}
    \bsrow{资本（净值）}{12}
    \bsrowtotal{合计}{110}
  }
\end{bsheet}

\begin{bsheet}
  \bsside{放贷后 · 同一家支行（万元）}{%
    \bsrow{准备金}{47}
    \bsrow{持有的债券}{63}
    \bsrow{对陈宁的贷款}{300}
    \bsrowtotal{资产合计}{410}
  }
  \bsside{负债 + 净值}{%
    \bsrow{客户存款}{98}
    \bsrow{陈宁的活期存款}{300}
    \bsrow{资本（净值）}{12}
    \bsrowtotal{合计}{410}
  }
\end{bsheet}

资产端多了一笔"对陈宁的贷款 300 万"，负债端同时多了一笔"陈宁的存款 300 万"。两行一起出现，账自动就平了。

王阿姨的余额那天晚上和前一天一模一样。没有任何一个储户的账户被动过一分钱。

而全社会的钱，那一秒多了 300 万。

\begin{pullquote}{}
银行不是把别人的钱借给你。银行是在给你贷款的同时，顺手把你的存款也造了出来。
\end{pullquote}

回到第 2 章的结论：你的存款说白了就是银行欠你的一张欠条。所以"银行创造货币"这件听着像魔法的事，拆开看很朴素，无非是给自己记了一笔"我欠陈宁 300 万"的账。

你我也能记。我现在就可以撕张纸写"欠陈宁 300 万"。

区别只有一条，但这一条是天堑：\textbf{银行写的欠条，全社会都当钱花。}陈宁拿它付给项老师，项老师拿它交契税、付中介费、给女儿报暑期班，没有一个人问过"这真的能兑现吗"。

那个没人问的瞬间，就是整台机器的地基。

故事往下走。项老师的卡不在这家银行。

跨行付款要在央行的账本上清算，这家支行必须真金白银划出 300 万\term{准备金}{reserves}\index{089@准备金}，也就是它存在央行账户里的那笔钱。可它账上只有 47 万。

去借。

\begin{egbox}{举个栗子：钱是先放出去，还是先凑够}
这家支行的操作顺序是：先在银行间市场拆入 253 万（准备金凑到 300 万，负债端多一笔"同业拆入 253 万"），然后把 300 万准备金划给项老师那家银行，陈宁的存款同时清零。

请盯住这个\textbf{时间顺序}。贷款早就放出去了，准备金是\hlmark{事后才去凑的}。银行不是先有准备金才敢放贷，是先放贷再补窟窿。窟窿补得贵不贵，看银行间市场当天的价，而那个价是央行说了算的。

\end{egbox}

\begin{bsheet}
  \bsside{陈宁跨行付款后（万元）}{%
    \bsrow{准备金}{0}
    \bsrow{持有的债券}{63}
    \bsrow{对陈宁的贷款}{300}
    \bsrowtotal{资产合计}{363}
  }
  \bsside{负债 + 净值}{%
    \bsrow{客户存款}{98}
    \bsrow{同业拆入}{253}
    \bsrow{资本（净值）}{12}
    \bsrowtotal{合计}{363}
  }
\end{bsheet}

贷款还在资产端，负债端却从"欠陈宁"变成了"欠同业"。

货币总量一分没少。那 300 万此刻躺在项老师的银行里，只是换了个门牌号。钱在银行之间流动，改变的是它记在谁名下，不是它存不存在。

能让它消失的只有一件事。

\section{陈宁每还一笔本金，就杀死一笔钱}

把时间快进到 2029 年。陈宁的公司被并购，他手上一笔期权兑现，决定把房贷一次性结清。那时候本金还剩 269 万。

银行的操作是：把他账上的 269 万存款划掉，同时把资产端"对陈宁的贷款"抹平。

\begin{bsheet}
  \bsside{结清之后（万元）}{%
    \bsrow{准备金}{47}
    \bsrow{持有的债券}{63}
    \bsrow{对陈宁的贷款}{0}
    \bsrowtotal{资产合计}{110}
  }
  \bsside{负债 + 净值}{%
    \bsrow{客户存款}{98}
    \bsrow{陈宁的存款}{0}
    \bsrow{资本（净值）}{12}
    \bsrowtotal{合计}{110}
  }
\end{bsheet}

那 269 万去哪儿了？

哪儿也没去。它不存在了。

它没进银行口袋。银行的资本还是 12 万，一分没多。这笔货币被注销了，像一条被抹掉的数据库记录。

陈宁每月还的那笔钱也一样。利息部分变成银行收入，会重新流回经济；\textbf{本金部分是被销毁的}，贷款和存款同步减少，社会上的钱真的少了那么一点。

\begin{databox}{这解释了正在发生的一件事}
几百万个陈宁同时决定"把房贷提前还了吧"，会怎么样？

在每一张家庭账本上这都是明智事：少欠钱，少付息，晚上睡得着。几百万张账本叠起来，它换了个名字，叫\textbf{货币紧缩}。存量贷款收缩，对应的存款被同步注销，社会总货币量往下走。

于是有了那个别扭的组合：央行说政策"适度宽松"，7 天期逆回购利率稳在 1.4\%（截至 2026 年 8 月），身边的人却普遍觉得钱很紧。央行能定水价，定不了水量。水量是千千万万个陈宁在"借"和"还"之间投票投出来的。

\end{databox}

\section{教科书里那个乘数，因果讲反了}

上过经济学原理课的话，脑子里大概有这么一套流程。

\begin{chainflow}
\chainnode{央行投放准备金}\chainarrow \chainnode{银行留下 10\% 法定准备金}\chainarrow \chainnode{把 90\% 贷出去}\chainarrow \chainnode{钱又存回银行}\chainarrow \chainnode{再贷出 81\%}\chainarrow \chainnode{最终放大 10 倍}
\end{chainflow}

这套东西叫\term{货币乘数}{money multiplier}\index{024@货币乘数}。它优雅、好考、能算出漂亮的等比数列。

它是错的。

不是"不够精确"，是因果方向整个反过来了。教科书说先有准备金才有贷款，现实是先有贷款再补准备金。

这不是什么非主流学派的私货。2014 年，英格兰银行在季度公报上发了一篇官方论文《现代经济中的货币创造》，白纸黑字写着：银行放贷并不是把储户的存款转手借出，而是通过放贷创造存款；并且明确指出，货币乘数式的叙述\hlmark{不是对货币实际创造过程的准确描述}。欧洲央行后来的官方文章、美联储研究人员的表述，口径完全一致。

十二年过去了。教科书还在照原样印。

\begin{mythbox}{常见误解}
\mvfrow{误解}{"存款准备金率决定货币总量。降准 0.5 个百分点，社会上的钱就成倍增加。"}
\mvfrow{事实}{准备金率今天是\textbf{成本和流动性的旋钮}，不是货币总量的乘法器。降准给银行释放长期低成本资金，负债成本下来一点，它\emph{可能}更愿意放贷。但如果没人想借，或者银行资本不够、风险偏好塌了，释放的钱就趴着当超额准备金，什么都不会发生。}
\end{mythbox}

\begin{mythbox}{常见误解}
\mvfrow{误解}{"银行是资金的搬运工，把储户的钱搬给借款人，中间赚个差价。"}
\mvfrow{事实}{银行是\textbf{货币的制造商}。它的生意是用自己的负债（存款）买别人的负债（贷款合同），并承担期限错配和信用风险。它之所以是被管得最严的行业，就因为它手里那支笔写出来的东西，全社会都当钱用。}
\end{mythbox}

\section{那什么在拦着银行}

既然敲两行字就能造钱，银行为什么不把全世界的钱都造出来，给自己发年终奖？

约束是真实存在的，只是不在你以为的地方。

\begin{flowlist}
  \flowitem{01}{有没有值得贷的人}{贷款合同要两个签名。银行想放，得有人敢借，还得让银行信他还得起。经济下行时最缺的从来不是钱，是\hlmark{敢签字的借款人}。}
  \flowitem{02}{资本充足率}{巴塞尔 III 的核心规则：每放一笔贷款，都要按风险权重占用一块自有资本。资本是股东真金白银掏的，不能凭空写两行。\textbf{这是最硬的一道闸。}}
  \flowitem{03}{赚不赚钱}{截至 2026 年二季度末，中国商业银行\term{净息差}{net interest margin}\index{029@净息差}只有 1.41\%，历史低位。放 100 万贷款一年毛赚 1.41 万，扣掉运营成本和坏账拨备，剩下的还未必比拆借给同业强。\hlmark{放贷不赚钱，机器就不想开。}}
  \flowitem{04}{监管在盯着}{房地产贷款集中度、行业限额、压力测试、逆周期资本缓冲。监管随时可以直接把某个方向的龙头拧小，不需要经过市场同意。}
  \flowitem{05}{央行定的那个价}{事后要补的准备金是有价格的。把这个价抬高，放贷成本上升，能过审的项目变少。\textbf{这才是央行真正的抓手。}}
\end{flowlist}

\begin{egbox}{举个栗子：资本这道闸有多硬}
某城商行核心一级资本 386 亿元。个人住房抵押贷款的风险权重按抵押率分档，这里取一个便于计算的档位 50\%，监管要求的核心一级资本充足率（含储备资本）按 7.5\% 算。

它能撑起的风险加权资产上限是 386 ÷ 7.5\% ≈ \textbf{5147 亿元}，按 50\% 权重折回去，房贷规模上限大约 \textbf{1.03 万亿元}。

落到陈宁头上：每放一笔 300 万的房贷，占用风险加权资产 150 万，占用核心一级资本 \textbf{11.25 万元}。这家银行最多放 34 万个陈宁，然后就得停。哪怕它趴在一座超额准备金的山上，也一笔都不能多放。

\end{egbox}

\section{三层货币：你以为的"钱"其实分三种电压}

现在装上这一章最有用的一副眼镜。

\begin{tablewrap}
\begin{xltabular}{\linewidth}{>{\raggedright\arraybackslash}p{0.20\linewidth}XXX}
\toprule
\thead{层} & \thead{是什么} & \thead{是谁的负债} & \thead{谁能拿着} \\
\midrule
\textbf{第一层} & 央行准备金 & 央行 & \textbf{只有银行}。银行之间清算专用，你我一辈子碰不到 \\
\textbf{第二层} & 商业银行存款 & 商业银行 & 你、我、企业、政府。日常说的"钱"几乎全在这一层 \\
\textbf{第三层} & 现金 & 央行 & 所有人。唯一能真攥在手里的那种 \\
\bottomrule
\end{xltabular}
\end{tablewrap}

拿电网打比方就清楚了。准备金是高压输电线上的电，只在变电站之间跑，你家电器插上去会炸。存款是入户的 220 伏，能烧水做饭。现金是充电宝里那点电，随身带着，量很小。

而把高压变成 220 伏的那台变压器，就是银行放贷。

这个比方能一路推下去：往高压线上加多少电压，你家插座都不会变一伏，除非变压器在工作。变压器停了，高压线上电压再高也点不亮一盏灯。第 8、9 章要讲的"传导失灵"，说的就是变压器不转。

\begin{pullquote}{}
第一层的钱和第二层的钱之间没有自动通道。这就是为什么央行"放水"，你的口袋可能一滴都感觉不到。
\end{pullquote}

2009 到 2019 年那十年，最能说明这副眼镜的用处。

发达国家搞了十年\term{量化宽松}{quantitative easing}\index{033@量化宽松}，央行资产负债表膨胀了好几倍。当时一大批人斩钉截铁地预言：恶性通胀马上就到，赶紧买黄金、囤罐头。

十年过去，CPI 通胀没来。

因为 QE 创造的是第一层。央行从银行手里买走国债，付给银行的是准备金。准备金躺在央行账户里，只能在银行之间转来转去，它不会自己长腿跑进陈宁的工资卡。

那些钱去哪儿了？它把国债和其他安全资产的收益率压到极低，逼着所有拿着钱的机构去买风险更高的东西：股票、房子、信用债、私募。所以那十年物价很平静，资产价格涨疯了。

存钱的人一无所获，持有资产的人富得流油。这不是道德问题，是管道问题。水注进了第一层的池子，而资产价格恰好是第一层水位最先溢出的地方。

要让钱真正进实体经济只有一条路：绕过银行的放贷决策，让钱直接落在居民和企业的存款账户上。这件事只有\textbf{财政}干得了。政府发工资、发补贴、给工程队结算，钱是直接打进第二层的，收款人的余额当场增加。

\begin{egbox}{举个栗子：同样是"放水"，结局差了四十年}
\textbf{2010 年代：QE。}美联储买国债，创造准备金（第一层）。结果是股票和房地产大涨，CPI 常年低于 2\% 的目标，美联储那些年发愁的是通胀太低。

\textbf{2020 到 2021 年：直接支票。}财政部把钱打进几乎每一户居民的账户（第二层），叠加失业救济加码、供应链断裂、劳动力紧张。结果是 2022 年 CPI 冲上四十年未见的高位。

同一个国家，同一批人，前后不到十年。差别不在"印了多少"，在\textbf{钱进了哪一层，以及拿到钱的人有没有立刻去买真东西}。

所以本书从不说"印钞必然通胀"。通胀是支出撞上供给产生的。2026 年通胀的粘性更多来自关税、去全球化和中东局势推高的能源成本，全是供给侧的事。

\end{egbox}

\section{稳定币是一家开在链上的窄银行}

把这一章的规律拿去套 2026 年最热闹的那个新东西，严丝合缝。

\term{稳定币}{stablecoin}\index{065@稳定币}怎么运作？你交 1 美元给发行方，发行方拿这 1 美元去买短期美国国债和货币市场工具，然后发给你一个代币，承诺随时按 1 比 1 换回美元。

翻译成这本书的语言：代币是发行方的负债，撑着它的是资产端的短期国债。它非常接近教科书里的\term{窄银行}{narrow bank}\index{082@窄银行}，只吸收资金、只投最安全的短期资产、不放贷，因此不创造新货币，只是把已有的钱换了条轨道。

截至 2026 年 8 月，全球稳定币市值在 3000 亿美元量级。美国的 GENIUS Act 已于 2025 年 7 月生效，为美元稳定币建立了联邦监管框架。监管者盯的问题和这本书一样：储备是什么，能不能随时兑付，发行方破产了谁赔。

\begin{warnbox}{小心}
"稳定"两个字是营销，不是物理定律。它稳不稳取决于资产端的质量和赎回机制，和一家银行会不会被挤兑是同一道题。历史上已经有算法稳定币在几天内归零的案例。

本书不对任何具体加密资产给投资建议，只提醒你用同一把尺子量它：\textbf{谁的负债？拿什么撑着？急着要钱的那天，那个"撑着"的东西卖得掉吗？}

\end{warnbox}

银行存款靠贷款和资本撑着，现金靠国家的征税能力撑着，稳定币靠短期国债撑着。任何"钱"，都是某个主体的负债。

\section{三百万亿里，只有 4\% 当过纸}

\begin{statrow}{3}
  \statitem{300+}{万亿}{中国 M2 规模（元，量级）}
  \statitem{10+}{万亿}{流通中现金 M0（元，量级）}
  \statitem{1.41}{\%}{商业银行净息差（2026 年二季度末）}
\end{statrow}

\begin{barfig}
  \barrow{银行存款（第二层，贷款创造）}{96}{约 96\%}
  \barrow{现金（第三层，央行印的）}{4}{约 4\%}
\end{barfig}

广义货币三百多万亿元，流通中的现金只有十几万亿元。95\% 以上的"钱"从没当过一天纸，它们是商业银行给某个陈宁、某家模具厂、某个地方平台批贷款时敲键盘敲出来的。

钱和债是同一枚硬币的两面。钱要变多就得有人多欠债；债要减少，钱必然变少。整套货币体系靠的是"持续有人愿意借"。

那如果哪天真的没人愿意借了呢？

日本试过。试了三十年，至今没试出一个能让别人照抄的答案。这件事目前没人知道解法。

\bookbreak

\begin{takeaway}{本章一句话}
绝大部分货币不是印出来的，是借出来的。银行放贷时同时创造资产和负债，还款时又把它注销掉。

\begin{itemize}
  \item 贷款创造存款。准备金是事后补的，不是原料。货币乘数那套讲法把因果讲反了。
  \item 真正的约束是借款人的需求、资本充足率、1.41\% 的净息差、监管和央行定的那个价，不是"手上有多少存款"。
  \item 钱分三层。QE 只灌第一层，所以它推资产价格胜过推物价；财政直接打钱灌第二层，所以 2021 年通胀来了。
  \item 碰到任何新的"钱"，先问：谁的负债，拿什么撑着。
\end{itemize}

\end{takeaway}

\begin{bridgebox}
\textbf{下一章}既然钱主要是商业银行造的，那个天天上新闻的央行到底在管什么？它管的不是水量，是水价和水温。
\end{bridgebox}

